\documentclass[10pt,letterpaper]{article}

\usepackage[letterpaper,margin=0.70in,headheight=22pt,headsep=16pt,footskip=24pt]{geometry}
\usepackage[T1]{fontenc}
\usepackage[utf8]{inputenc}
\usepackage{lmodern}
\usepackage{microtype}
\usepackage{xcolor}
\usepackage{graphicx}
\usepackage{booktabs}
\usepackage{array}
\usepackage{tabularx}
\usepackage{longtable}
\usepackage{amsmath,amssymb}
\usepackage{tikz}
\usetikzlibrary{arrows.meta,positioning,calc,shapes.geometric,decorations.pathreplacing,fit}
\usepackage{enumitem}
\usepackage{fancyhdr}
\usepackage{titlesec}
\usepackage{caption}
\usepackage{listings}
\usepackage{xurl}
\usepackage[hidelinks]{hyperref}
\usepackage{lastpage}
\usepackage{float}
\usepackage{placeins}

\definecolor{TrawlAmber}{HTML}{A66A00}
\definecolor{TrawlGray}{HTML}{555555}
\definecolor{TrawlBlue}{HTML}{144E8C}
\definecolor{LightGray}{HTML}{F3F3F3}
\definecolor{MidGray}{HTML}{D9D9D9}
\definecolor{DarkText}{HTML}{1A1A1A}
\definecolor{GoodGreen}{HTML}{2B7A2B}
\definecolor{WarnOrange}{HTML}{A65E00}

\renewcommand{\familydefault}{\sfdefault}
\setlength{\parindent}{0pt}
\setlength{\parskip}{5pt}
\renewcommand{\arraystretch}{1.15}
\sloppy
\setlist[itemize]{leftmargin=1.2em,itemsep=2pt,topsep=2pt}
\setlist[enumerate]{leftmargin=1.4em,itemsep=2pt,topsep=2pt}

\titleformat{\section}{\Large\bfseries\color{DarkText}}{}{0pt}{}[\vspace{-3pt}\color{TrawlAmber}\titlerule]
\titleformat{\subsection}{\large\bfseries\color{DarkText}}{}{0pt}{}
\titleformat{\subsubsection}{\normalsize\bfseries\color{DarkText}}{}{0pt}{}

\newcommand{\code}[1]{\texttt{\detokenize{#1}}}
\newcolumntype{Y}{>{\raggedright\arraybackslash}X}
\newcolumntype{L}[1]{>{\raggedright\arraybackslash}p{#1}}
\captionsetup{font=small,labelfont=bf}

\lstdefinestyle{shell}{
  basicstyle=\ttfamily\footnotesize,
  breaklines=true,
  columns=fullflexible,
  frame=single,
  rulecolor=\color{MidGray},
  backgroundcolor=\color{LightGray},
  keepspaces=true,
  showstringspaces=false
}
\lstset{style=shell}

\newcommand{\notebox}[1]{%
  \vspace{4pt}\noindent\fcolorbox{TrawlBlue}{LightGray}{%
  \begin{minipage}{0.965\linewidth}%
  \textbf{Note:} #1%
  \end{minipage}}\vspace{4pt}}

\newcommand{\importantbox}[1]{%
  \vspace{4pt}\noindent\fcolorbox{TrawlAmber}{LightGray}{%
  \begin{minipage}{0.965\linewidth}%
  \textbf{Important:} #1%
  \end{minipage}}\vspace{4pt}}

\newcommand{\field}[1]{\texttt{\detokenize{#1}}}

\pagestyle{fancy}
\fancyhf{}
\fancyhead[L]{\small Datatrawl-UG001}
\fancyhead[C]{\small Storage-Safe Archive Trawling Engine v0.2.0}
\fancyhead[R]{\small User Guide}
\fancyfoot[L]{\small Datatrawl-UG001 v1.1}
\fancyfoot[C]{\small West Virginia University}
\fancyfoot[R]{\small Page \thepage\ of \pageref{LastPage}}
\renewcommand{\headrulewidth}{0.8pt}
\renewcommand{\headrule}{\hbox to\headwidth{\color{TrawlAmber}\leaders\hrule height \headrulewidth\hfill}}
\renewcommand{\footrulewidth}{0.8pt}
\renewcommand{\footrule}{\hbox to\headwidth{\color{TrawlAmber}\leaders\hrule height \footrulewidth\hfill}}

\title{datatrawl Storage-Safe Archive Trawling Engine v0.2.0\\User Guide}
\author{West Virginia University}
\date{Datatrawl-UG001 v1.1 \\ July 3, 2026}

\begin{document}
\thispagestyle{fancy}
{\sffamily
\begin{minipage}[t]{0.39\textwidth}
  {\fontsize{24}{26}\selectfont\bfseries\color{TrawlAmber}datatrawl}\par
  {\small Datatrawl-UG001 v1.1, July 3, 2026}
\end{minipage}%
\begin{minipage}[t]{0.61\textwidth}
  \raggedleft
  {\large\bfseries User Guide}\par
  {\small Install, plan, run, resume.}\par
  {\small West Virginia University Radio Astronomy Instrumentation Lab (RAIL)}
\end{minipage}
}
\vspace{6pt}
{\color{TrawlAmber}\hrule height 1.2pt}
\vspace{8pt}

\section{The Model in One Paragraph}
Every datatrawl run is four choices: a \textbf{telescope} (geometry and
metadata), a \textbf{source} (where files come from), a \textbf{reader} (how
one file becomes arrays), and an \textbf{analyzer} (the science). The engine
then streams the archive one file at a time --- fetch, analyze, delete,
checkpoint --- into a single resumable product per run. \code{datatrawl
doctor} exists to make the four choices and their prerequisites legible
\emph{before} you spend hours on a scan.

\section{Installation}
\begin{lstlisting}
git clone https://github.com/WVURAIL/datatrawl.git
cd datatrawl
python -m venv .venv && source .venv/bin/activate
pip install -e ".[dev]"
datatrawl --version
pytest -q
\end{lstlisting}
Optional extras: \code{[cadc]} for CADC fetch, \code{[survey]} for Datatrail
inventory builds, \code{[gpu]} for GPU analyzers (CuPy is installed separately
--- see \S\ref{sec:cupy}).

\section{Discovering What Is Available}
\begin{lstlisting}
datatrawl list                 # everything at a glance
datatrawl list telescopes      # ready vs. stub
datatrawl list readers
datatrawl list analyzers
datatrawl list sources
\end{lstlisting}
Then verify one concrete combination end to end:
\begin{lstlisting}
datatrawl doctor --telescope chime --source cadc-datatrail \
                 --reader chime-baseband --analyzer spectrum
\end{lstlisting}
\code{doctor} reports, per choice, what is installed, what is configured, what
credentials are missing, and whether the combination is ready to run.

Before committing to any of that, map the archive itself. Recon is
\code{survey}'s no-download mode: it lists datasets across scopes and writes
the result as a reusable map (\code{scopes.jsonl}, or
\code{scopes-<name>.jsonl} with \code{--name}, so successive recons do not
overwrite each other):
\begin{lstlisting}
datatrawl survey --telescope gbo --scopes-only              # this telescope's scopes
datatrawl survey --scopes-only                              # every scope datatrail sees
datatrawl survey --telescope gbo --scopes-only \
    --scope gbo.acquisition.processed --match gain --expand --name gains
\end{lstlisting}
\code{--telescope} narrows the walk to that telescope's \emph{live} scopes
(first name component --- deliberately not the instrument YAML, which declares
only what the event survey walks); omit it for zero-knowledge discovery.
\code{--match} keeps name matches (comma-separated terms, ANDed);
\code{--expand} opens each kept dataset one level, so a container such as
\code{complex_gains} yields its children as rows (with a \field{parent}
column). Every row resolves directly with
\code{datatrail ps <scope> <dataset> -s}. A scope datatrail fails to answer
for is named in-walk and in a closing \code{[warn]} block marking the map
INCOMPLETE --- an outage is never mistaken for emptiness.
\notebox{Recon finds names; survey's event walk enumerates \emph{event-keyed}
data only. A container of timestamped, non-event products (calibration
acquisitions) is reached through recon \code{--expand} and the programmatic
file listing \code{DATATRAIL.files()}, never through an event survey. The
split is spelled out in \code{docs/DATATRAIL_BOUNDARY.md}.}

\section{Building an Inventory}
\code{survey} walks the archive once and writes a named inventory directory
(default \code{data/<telescope>-fid<freq_ids>/}, or \code{--name}); it records
the telescope, source, and reader choices inside the inventory so later scans
need only the analyzer:
\begin{lstlisting}
datatrawl explore --telescope chime --source cadc-datatrail   # look, no download
datatrawl survey  --telescope chime --source cadc-datatrail   # -> data/chime.../
\end{lstlisting}
Inventories stream from disk during a scan; the engine never needs the whole
archive listing in memory.

The reader also owns the \emph{archive file shape}: which files one event
contributes and what they are named (\field{Reader.survey_files}). Survey
resolves the telescope's canonical reader (override with \code{--reader}) and
writes each verified file's \field{name} into its row, so a later scan stages
exactly what the survey verified --- the two cannot drift. A non-baseband
product becomes surveyable by loading a small external shape reader with
\code{--plugin}; see \code{docs/ADDING_A_READER.md}.

\section{Running a Scan}
\begin{lstlisting}
datatrawl scan --name chime --analyzer spectrum --select 614,706
\end{lstlisting}
\begin{itemize}
  \item \code{--select} is interpreted by the analyzer. For the spectrum
        analyzer: explicit \field{freq_id}s such as \code{844},
        \code{614,706}, or ranges like \code{506-552}. Event-oriented
        analyzers take \code{events:349382977[,...]}, and a
        \field{plan_runs} may return the dict form
        \code{{"events": [...], "freq_ids": ...}}; filters are exact and
        ANDed, and a malformed selection fails with the grammar named, not a
        traceback. A multi-item selection becomes independent products that
        resume and fail independently.
  \item \code{--checkpoint-every N} (default 50) sets how many files may need
        re-processing after a kill; each checkpoint is an atomic full-product
        save.
  \item \code{--download-workers} / \code{--max-staged-files} parallelize
        \emph{fetch only}; the engine refuses values above 1 for analyzers
        that declare \field{requires_in_order}.
  \item \code{--plugin yourpkg.module} loads an external plugin explicitly;
        installed packages advertising the \code{datatrawl.plugins}
        entry-point group are found automatically.
\end{itemize}

\importantbox{Storage safety is the point: at the defaults, at most one staged
file exists at a time, it lives in a per-invocation staging directory
(concurrent scans cannot delete each other's files), and it is deleted as soon
as the analyzer has consumed it. Do not point the staging area at precious
data.}

\section{Kill and Resume}
A scan may be killed at any instant --- session timeout, preemption,
\code{Ctrl-C}. On restart with the same configuration and output, the analyzer
re-loads its product, skips every unit already accumulated, and continues. The
guarantee: at most the files consumed since the last checkpoint are re-done,
and the product is never left partially written (saves go to \code{.tmp}, then
rename). A resume against an \emph{incompatible} existing product is refused
with an explicit error rather than silently mixing two configurations; use a
clean output location to rebuild.

\section{CANFAR / CADC Notes}
On a CANFAR session: install with the \code{[cadc]} (and, for inventory
builds, \code{[survey]}) extras, ensure a valid CADC proxy certificate, and
let \code{datatrawl doctor --source cadc-datatrail ...} confirm credentials
and the Datatrail contract before launching. Datatrail remains the archive
authority --- datatrawl only reads; the boundary is specified in
\code{docs/DATATRAIL_BOUNDARY.md}. For long scans, size
\code{--checkpoint-every} so an unexpected session end costs minutes, not
hours.

\section{GPU Sessions and CuPy}\label{sec:cupy}
GPU analyzers use CuPy, which must match the session image's CUDA runtime, so
it is not pinned as a dependency. On a GPU node:
\begin{lstlisting}
datatrawl setup-cupy            # report the image's CUDA + the matching wheel
datatrawl setup-cupy --install  # install cupy-cudaXXx into the active venv
\end{lstlisting}
CPU-only environments need nothing; the built-in analyzers fall back to
NumPy with identical arithmetic.

\section{Writing an Analyzer}
The analyzer owns its product file. Implement the lifecycle and register it:
\begin{lstlisting}
class MyAnalyzer(Analyzer):
    info = PluginInfo(name="my-analyzer", ...)
    requires_in_order = False   # True for any order-dependent product

    def resume(self, path, ctx):      # load existing product; raise if
        ...                           # present-but-incompatible
    def processed_keys(self):         # unit keys already in the product
    def begin(self, ctx, first_meta): # once, at the first new file
    def consume_file(self, arrays, meta):  # accumulate; return n frames
    def save(self, path):             # atomic write with provenance
    def summary(self):                # one status line
\end{lstlisting}
Two optional hooks shape selection: \field{resolve_selection} interprets the
user's \code{--select} for your analysis, and \field{plan_runs} splits it into
independent runs, each with a fresh analyzer instance and its own resumable
product. All accumulator writes happen on the engine's main thread, so no
locking is needed. The complete walkthrough, including the inventory-metadata
contract your reader must satisfy, is in \code{docs/ADDING_AN_ANALYZER.md}
(and its \code{READER}/\code{SOURCE}/\code{TELESCOPE} siblings).

Two patterns built on those hooks are worked end to end in the repository.
\emph{Per-event fan-out}: \field{plan_runs} returns one
\code{{"events": [ev], "freq_ids": spec}} sub-selection per event, so each
event becomes its own resumable product
(\code{ev<event>[_<freq_ids>].npz}); both archive and local sources
understand the dict natively. \emph{Companion side-loads}: an analyzer
needing a per-event auxiliary (calibration gains, flags) loads its lookup
once in \field{begin} and resolves per file off \field{meta["event"]},
stamping the companion's identity into the product and refusing a resume if
the lookup later disagrees. Day-keyed companion archives resolve lazily with
\code{DATATRAIL.files(scope, day)}, imported from
\code{datatrawl.plugins.sources} --- the programmatic
\code{datatrail ps -s}. The runnable, CLI-tested reference is
\code{examples/per_event_companions.py}; the patterns are written out in
\code{docs/ADDING_AN_ANALYZER.md}.

\importantbox{If your product depends on the order files are consumed --- a
running or trailing statistic, a CFAR baseline --- set
\field{requires_in_order = True}. The engine will then refuse the parallel
settings that deliver files in download-completion order, turning a silent
wrong answer into a loud configuration error.}

\section{Troubleshooting}
\begin{tabularx}{\linewidth}{L{0.34\linewidth} Y}
\toprule
\textbf{Symptom} & \textbf{First move} \\
\midrule
``How do I even start?'' & \code{datatrawl list}, then \code{datatrawl doctor}
with your four choices. \\
Scan refuses to start & Read the doctor-style preflight error: missing extra,
missing credential, or a reader/telescope mismatch. \\
Resume refused & The existing product was built with different parameters;
that refusal is the safety feature. Rebuild in a clean output location. \\
Parallel flags refused & The analyzer declares \field{requires_in_order};
run at the defaults. \\
Anything else & \code{docs/TROUBLESHOOTING.md} in the repository. \\
\bottomrule
\end{tabularx}

\section{Related Documents}
\textbf{Datatrawl-DS001} (product specification: contracts, limits,
interfaces); the \code{docs/ADDING_*.md} extension guides;
\code{docs/DATATRAIL_BOUNDARY.md}; and, for the RAIL detector built on this
engine, \textbf{PilotProxy-DS001/UG001}.

\section{Revision History}
\begin{tabularx}{\linewidth}{L{0.12\linewidth} L{0.20\linewidth} Y}
\toprule
\textbf{Rev} & \textbf{Date} & \textbf{Changes} \\
\midrule
v1.0 & July 1, 2026 & Initial release, documenting engine v0.1.0. \\
v1.1 & July 3, 2026 & Engine v0.2.0: recon (\code{--scopes-only} with
\code{--match}/\code{--expand}/\code{--name}, telescope-scoped), event
selection grammar, reader-owned archive file shape (\code{survey --reader}),
\code{DATATRAIL.files()}, per-event fan-out and companion patterns. \\
\bottomrule
\end{tabularx}

\end{document}
